%! Tex main = MarquezSalazarBrandon-PropuestaDeTesis.tex

\section{Cronograma}

% --- Lógica de automatización ---
\newcounter{weekstart}
\newcounter{weekend}
\newcounter{lastmaxend}

\newcommand{\ganttparallel}[4]{%
    \setcounter{weekstart}{#2}%
    \setcounter{weekend}{\value{weekstart} + #3 - 1}%
    \ganttbar[bar/.append style={fill=#4}]{#1}{\value{weekstart}}{\value{weekend}}%
    % Guardamos el final más lejano por si queremos retomar la serie después
    \ifnum\value{weekend}>\value{lastmaxend}
        \setcounter{lastmaxend}{\value{weekend}}
    \fi
}

\newcommand{\ganttseries}[3]{%
    \setcounter{weekstart}{\value{lastmaxend} + 1}%
    \setcounter{weekend}{\value{weekstart} + #2 - 1}%
    \ganttbar[bar/.append style={fill=#3}]{#1}{\value{weekstart}}{\value{weekend}}%
    \setcounter{lastmaxend}{\value{weekend}}
}

% --- En tu documento ---
\begin{center}\resizebox{\textwidth}{!}{%
    \hspace*{-1.2cm}
    \begin{ganttchart}[
        hgrid, vgrid,
        x unit=0.2cm, % Ajusta el ancho total
        y unit title=0.7cm,
        y unit chart=0.6cm,
        title/.style={fill=none, draw=none, font=\small\bfseries},
        bar/.style={draw=black, fill=blue!30},
        canvas/.style={draw=black, fill=none},
        milestone/.style={fill=black},
        progress label text={}, % Limpia etiquetas de porcentaje
        group right shift=0,
        group top shift=0.7,
        group height=.3,
        group peaks tip position=0,
        bar label node/.style={
            anchor=east, 
            text width=5cm,  % <--- Ajusta este valor según necesites
            align=right,     % <--- Necesario para que respete el ancho
            font=\footnotesize      % Opcional: reducir fuente si el texto es muy largo
        },
        milestone label node/.style={
            anchor=east, 
            text width=5cm,  % <--- Ajusta este valor según necesites
            align=right,     % <--- Necesario para que respete el ancho
            font=\footnotesize\itshape      % Opcional: reducir fuente si el texto es muy largo
        },
    ]{1}{52}
            
        % Título de Meses (4 o 5 semanas por mes)
        \gantttitle{Cronograma Anual 2026 (Semanas)}{52} \\
        \gantttitle{Ene}{4} \gantttitle{Feb}{4} \gantttitle{Mar}{5} 
        \gantttitle{Abr}{4} \gantttitle{May}{4} \gantttitle{Jun}{5}
        \gantttitle{Jul}{4} \gantttitle{Ago}{4} \gantttitle{Sep}{5}
        \gantttitle{Oct}{4} \gantttitle{Nov}{4} \gantttitle{Dic}{5} \\
        
        % Inicializar el contador en 0 para empezar en enero (Mes 1)
        \setcounter{lastmaxend}{0}
        \setcounter{weekend}{0}

        % --- Definición de Etapas ---
        \ganttparallel{Revisión de la literatura}{1}{4}{Violet!20} \\
        \ganttparallel{Investigación sobre \textit{Framewoks}}{1}{2}{Maroon!20} \\
        \ganttparallel{Implementación y validación de la función objetivo}{3}{3}{PineGreen!20} \\
        \ganttseries{Pruebas con SGA}{4}{Orange!20} \\
        \ganttseries{Pruebas con PSO}{4}{Orange!20} \\
        \ganttseries{Pruebas con EDAs}{6}{Orange!20} \\
        \ganttseries{Pruebas con ACO}{4}{Orange!20} \\
        \ganttseries{Pruebas con IWO}{4}{Orange!20} \\
        \ganttseries{Pruebas con VCS}{4}{Orange!20} \\
        \ganttseries{Pruebas con ASO}{4}{Orange!20} \\
        \ganttseries{Análisis de resultados}{3}{Black!20} \\
        \ganttparallel{Redacción de la Tesis}{27}{20}{Red!80} \\

        % Hitos (Milestones)
        \ganttmilestone{Primer reporte de avances}{11} \\
        \ganttmilestone{Segundo reporte de avances}{23} \\
        \ganttmilestone{Defensa de Tesis}{49}

    \end{ganttchart}

    \vspace{10pt}
}
    \pgfornament[width=3cm]{88} % El toque final
\end{center}
